\documentclass[title=Monsieur, name=Médiocre, r=0.75, g=0.49, b=0.89]{../mrms}
\begin{document}

\coverpage{../img/barnak.jpg}
\collectionpage
\titlepage
\txtpage{
    Monsieur Médiocre avait un secret. Il était en fait un dieu déchu.

    Il venait d'une montagne où résident toutes sortes de dieux. Certains l'appelent le mont Olympe,
    d'autres l'appeleraient La Grande Montagne\texttrademark.

    Mais monsieur Médiocre avait toujours eu des idées de grandeurs...
}
\txtpage{<Montagne\texttrademark>}
\txtpage{
    Un jour, il décida d'aller voir Drareg, un très, très grand dieu.

    - Désormais, je serai le dieu des dieux! Déclara monsieur Médiocre.

    Drareg se mit alors à rire, car monsieur Médiocre était le dieu des mauvaises idées et des délusions
    de grandeur.
}
\txtpage{<Drareg (on voit juste les pieds)>}
\txtpage{
    Drareg riait si fort, que son souffle emporta monsieur Médiocre, et le fit dégringoler tout en
    bas de la montagne. Monsieur Médiocre était ainsi banni de La Grande Montagne\texttrademark, et
    avait perdu son statut de dieu.
}
\txtpage{<Barnak qui dégringole>}
\txtpage{
    Monsieur Médiocre se dit que s'il ne pouvait pas être le dieu des dieux, alors il pourrait très
    bien être le dieu des monsieurs et des madames.

    Il commença par aller voir la première personne qu'il pu trouver.

    - Prosternes-toi devant ton nouveau dieu! Annonça monsieur Médiocre.

    - Boop? Boop boop! Répondit monsieur Lette.

    Confus, monsieur Médiocre partit chercher une deuxième personne.
}
\txtpage{<barnak avec cossin lette>}
\txtpage{
    - Désormais, tu me vénéreras, je suis ton dieu! Exigea monsieur Médiocre à la prochaine personne.

    - Quoi? Vous êtes qui? Vous n'avez même pas de... de... commença madame Fin-De-Phrase. Enfin
    vous savez... un grand bâtiment qui sert à défendre des envahisseurs! Et visiblement vous n'avez
    pas de... enfin je vais trouver le mot... ah oui! De reine.

    - Huh? Répondit monsieur Médiocre. Vous voulez dire un château? J'ai besoin d'une reine et d'un
    château?

    - J'imagine. Les rois adorent dirent qu'ils sont des dieux, et ils ont des châteaux et des
    reines, finit madame Fin-De-Phrase.
}
\txtpage{<barnak avec mme fin-de-phrase>}
\txtpage{
    Un peu confus, monsieur Médiocre se mit donc à la recherche d'un entrepreneur général qui pourrait
    lui construire un château digne de ce nom. Ce serait une première étape. Il avait trois critères:
    Que ce soit bon, que ce soit beau, et que ce soit pas cher.
}
\txtpage{<barnak qui rêve à un bô chateaux}
\txtpage{
    Après avoir rencontré plusieurs entrepreneurs qui étaient tous trop dispendieux, monsieur
    Médiocre finit par rencontrer monsieur Ça-Fait-La-Job, qui lui expliqua:

    - Ah bah oui j'peux vous faire ça ben vite, et vous verrez, on pourra s'en sortir pour \textit{cheap},
    Ayez confiance monsieur Médiocre, ce sera à votre hauteur!

    Extraordinaire! Monsieur Médiocre allait enfin pouvoir avoir le château qu'il méritait.
}
\txtpage{<m. médiocre avec m ça-fait-la-job>}
\txtpage{
    Monsieur Ça-Fait-La-Job se mit au travail. Première étape: Faire des plans! D'habitude un
    château ça a un donjon, et des remparts. Il dessina donc un carré, puis une ligne qui faisait le
    tour.

    Monsieur Médiocre regarda le plan.

    - C'est assez sommaire comme plan. Vous croyez que ce sera suffisant pour inspirer le respect et
    arrêter de vilains preux?

    - Oh vous savez, répondit monsieur Ça-Fait-La-Job, je peux vous dire que ça fera la job.

    Monsieur Médiocre lui donna donc le go pour commencer les travaux.
}
\txtpage{<le plan>}
\txtpage{
    D'abord, monsieur Ça-Fait-La-Job construisit le rempart.

    - Ça ressemble à quoi notre stock de pierres? Demanda-t-il à monsieur Médiocre.

    - Oh... il semble y avoir des pierres un peu partout sur le terrain, est-ce que ce serait suffisant?
    Lui répondit monsieur Médiocre.

    - Ça fera la job, conclut monsieur Ça-Fait-La-Job.
}
\txtpage{<champ avec des pierres>}
\txtpage{
    Après un long moment à collecter toutes les pierres, monsieur Ça-Fait-La-Job ne réussit à construire
    que deux couches avant d'être complètement à court.

    Il opta donc pour terminer en utilisant les lattes de cloture en bois qu'il avait déjà.
}
\txtpage{<rampart (c'est une cloture)>}
\txtpage{
    Monsieur Ça-Fait-La-Job tourna ensuite son attention vers le château. N'ayant déjà plus de pierres,
    il utilisa donc le bois restant de sa dernière construction.

    Afin que le château soit le plus haut possible, il commença par quatres pilliers, avec une plateforme dessus.
    En haut de cette plateforme, il construisit la place forte du château: Quatres bons murs, et un toit.

    Pour garantir l'accessibilité, il installa une rampe, avec un angle respectable de vingt degrés.
}
\txtpage{<château (c'est un parc)>}
\txtpage{
    Fier de son travail, monsieur Ça-Fait-La-Job alla chercher monsieur Médiocre pour le lui présenter.

    - Ce fut rapide, s'exclama monsieur Médiocre! Bon, ce n'est peut-être pas aussi grandiose que je l'espérais,
    mais si vous me dites que c'est bon...

    Et monsieur Ça-Fait-La-Job hocha la tête de façon nonchalante.
}
\txtpage{<m médiocre et m ça fait la job avec le "château">}
\txtpage{
    Avant de pouvoir organiser sa cérémonie de déification, monsieur Médiocre devait alors trouver
    une reine.

    Il s'allia avec un groupe ultra-méga-secret de ninjas, qui lui promirent que la meilleure solution
    était de capturer la bonne personne. Ils l'assurèrent qu'ils connaissaient la madame qui serait
    parfaite comme reine.

    Monsieur Médiocre leur fit grâce de son aide directe. Hélas, ces satanés ninjas échouèrent lamentablement...
    Mais ça c'est une autre histoire.
}
\txtpage{<M médiocre avec un ninja>}
\txtpage{
    Un peu plus tard, monsieur Médiocre ruminait dans son château. Il devait trouver des nouvelles idées
    pour obtenir une reine.

    C'est alors que sa porte s'ouvrit brusquement.
}
\txtpage{<Monsieur médiocre qui rumine>}
\txtpage{
    Dans la porte se tenait monsieur Preux. On le sait car il était coiffé du haume caractéstique d'un preux chevalier.
    Monsieur Preux était venu pour arrêter les plans machiavéliques de monsieur Médiocre. Il avait même
    apporté l'épée légendaire: Excalibarre.

    Monsieur Médiocre dégaina son épée.
}
\txtpage{<Monsieur Preux>}
\txtpage{
    Puisque monsieur Preux ne disait rien, monsieur Médiocre prit la parole:

    - C'est une bonne chose tout compte fait. Minable chevalier! Bats-toi!

    - Boop, boop! Boop! (Non, Barnak! C'est à toi de battre! Dieu médiocre! Bats-toi!) Rétorqua monsieur Preux.
}
\txtpage{<Monsieur médiocre qui monologue>}
\txtpage{
    Outragé par l'impolitesse d'avoir utilisé son prénom plutôt que son nom de famille, monsieur
    Médiocre brandit son épée, et s'avança vers monsieur Preux...
}
\txtpage{<Monsieur médiocre qui attaque (m preux en contrejour?)>}
\txtpage{
    Son pied cogna un clou mal rentré, il trébucha, et il tomba par terre!
}
\legalpage{13 septembre 2026}
\backpage

\end{document}
